\documentclass[fontsize=16pt]{mybib}
\author{OTS}
\renewcommand{\DFA}{12mm}
\renewcommand{\DZA}{1.3}

\graphicspath{{../../../assets/images/}}
\setheaderlogo{mnr-green.png}

% Absatzformatierung
\setlength{\parindent}{0pt}
\setlength{\parskip}{0.6em}
\renewcommand{\arraystretch}{0.8}

\title{Nehemia 1}
\author{Lothar Schmid}
\date{}

\begin{document}
\mytitlepage
\tableofcontents
\flushbottom
\newpage

\section{Teil 1: Nehemia erörtert das Problem 1 - 4a}

\biblezit{Neh}{1:1}

Es war im Herbst, zwischen November und Mitte Dezember, um 444 \vChr. Alles fing auf einer Burg oder Palast in Susa an. \enquote{Artahsasta Longimanus} oder auch Artaxerxes genannt, war zu der Zeit König von Persien. Susa liegt im Westen des heutigen Iran neben dem Zagros-Gebirge. Dieser Palast ist heute archäologisch gut erforscht. Es ist der gleiche Palast, in dem Esther gewirkt hat und Daniel in \biblezit{Dan}{8:2} seine Vision hatte.

Hier in dieser Burg war ein junger Mann als Mundschenk des Königs, mit dem Namen Nehemia, tätig. Der Mundschenk war nicht nur einfach ein Kellner, ein Mundschenk war ein Vertrauter des Königs. Dieser Posten war mit einer grossen Verantwortung und grossen Vorrechten verbunden. Bei Mahlzeiten probierte er den Wein, um sicherzustellen, dass dieser nicht vergiftet war. Der Mundschenk musste gut aussehend und gebildet sein und sollte auch ein fröhliches Gemüt haben. Könnt ihr euch vorstellen, wie das ist? Wenn euer Vorkoster, an dem euer Leben hängt, missmutig und schlecht gelaunt erscheint, und ihr nicht wisst, ob er etwas Schlechtes getrunken hat? Würdet ihr euch da noch sicher fühlen? 
So ist es auch nicht verwunderlich, dass dem König die Traurigkeit seines Mundschenks sofort aufgefallen ist. 

Das grössere Problem für den Mundschenk, waren die Launen des Königs, diese konnten einem Mundschenk ernste Probleme bereiten, sogar bis zur Gefahr mit einem Kopf kürzer in den Feierabend zu gehen. Darum  waren gute Laune und Optimismus in dem Job wichtig. In \biblezit{IMos}{40:1} lesen wir eine Geschichte über einen Mundschenk, der fast hingerichtet wurde. Aber wieder zurück auf unsere Burg.

Nehemia war der Mundschenk von König Artaxerxes I. der von 465 bis 424 \vChr{} regierte. Nehemia war der Sohn Hachaljas und er war ein Jude. Da es zu der Zeit keine Nachnamen gab, war es üblich, den Namen des Vaters mit anzugeben, um sich von Gleichnamigen zu unterscheiden. \biblezit{Neh}{3:16}, \biblezit{Esr}{2:2} Nehemia gehörte zu den Nachkommen der Juden, die ein paar Generationen vorher von den Babyloniern aus Israel verschleppt und vertrieben wurden \biblezit{IIChr}{36:20}.

Es war nicht ungewöhnlich, dass ein Jude einen so verantwortungsvollen Posten im Exil bekam. Wir lesen von mehreren Juden in der Bibel, die eine hohe Stellung einnahmen, so der Prophet Daniel und seine Freunde, Mordechai und seine Nichte Esther. Alle zeichneten sich durch ihre hohe Integrität und Loyalität gegenüber ihren Vorgesetzten aus.

Das ist auch wichtig für uns. Wir Christen sollten Menschen sein, auf die sich der Arbeitgeber, der Vermieter und auch die Obrigkeit verlassen kann. Wir sollten immer ehrlich gegenüber unseren Vorgesetzten sein. Die Menschen sollten uns vertrauen können. Wir sind ehrlich, nicht neidisch und suchen keinen Vorteil auf Kosten unserer Mitmenschen. Oft hört man: \enquote{Was du bist Christ? Was du da machst, ist aber gar nicht christlich.} Die Menschen schauen auf dich und mich, besonders wenn sie wissen, dass wir Christus nachfolgen.

\biblezit{Neh}{1:2}\\
Ein Teil der verschleppten Juden sind schon vor knapp 100 Jahren zurück nach Israel gegangen. Obwohl Nehemia im Exil geboren wurde, war er im Herzen immer noch Jude und hatte ein Herz für sein Volk.  Als Hanani, der Bruder von Nehemia, zusammen mit ein paar Männern aus Jerusalem zurückgekehrte, ging er zu seinem Bruder und erkundigte sich nach der Lage, die aktuell in Jerusalem herrschte.

\biblezit{Neh}{1:3}\\
Der Bericht war katastrophal, wie wir in Vers 3 lesen. Jerusalem, das Mittelpunkt und Hauptstadt des Landes, die Stätte Gottes, lag offen da. Die Mauer waren niedergerissen und die Tore verbrannt.

Stadtmauern waren nicht nur zur Sicherheit da. Für uns heute haben Mauern und Stadttore keine grosse Bedeutung mehr, aber zu der damaligen Zeit zeigten die Mauern, Stärke, Macht, Unabhängigkeit und sie gaben Schutz vor Feinden und wilden Tieren. Je grösser und dicker die Mauer war, umso mächtiger vermutete man ein Herrscher dahinter.

\section{Teil 2: Nehemia betet um Hilfe 4b - 11}

\biblezit{Neh}{1:4}

Als Nehemia die Botschaft vernommen hatte, weinte er, fastete und betete zu Gott.

Wie ist das mit uns? Wenn wir im Fernsehen hören, dass in Nigeria wieder Christen ermordet wurden, nur weil sie an Christus festhielten? Sind wir noch erschüttert? Oder lässt uns das kalt? Unter den vielen Nachrichten die wir heute über das Fernsehen, Internet, Zeitungen und Handy alle paar Minuten bekommen, werden wir von schlechten Nachrichten überflutet. Aber nicht nur das, es wird auch kaum über die Christenverfolgung berichtet.

Weil die öffentlichen Medien wenig über verfolgte Christen sprechen, umso wichtiger ist es, dass wir uns selber diese Informationen beschaffen. So wie auch Nehemia sich erkundigt hat können auch wir uns unsere Information einholen. Zum Beispiel bei Open Doors oder HMK. Diese Organisationen arbeiten mit verfolgten Christen zusammen, helfen und unterstützen sie. Alle Christen, die Jesus angenommen haben, gehören zum gleichen Leib, bei dem das Haupt der Christus ist. So schreibt Paulus in \biblezit{IKor}{12:26}.

\biblezit{Neh}{1:5-11}

Nehemia weinte und fastete nicht nur, sondern er betete auch jeden Tag. Interessant ist, dass in Vers 4 steht: \enquote{etliche Tage lang}. Das heisst, er hat nicht nur einmal gebetet und Gott hat ihn erhört, sondern es ging länger. Wenn wir diesen Bibeltext so lesen, könnten wir denken, dass alles schnell ging. Aber so war es nicht. \enquote{Etliche Tage lang} heisst es dort. Beten ist eine Fleissarbeit.

Wir kennen den Ausgang eines Gebetes nicht. Georg Müller, der Mann, der in England Waisenhäuser für mittellose Kinder aufgebaut hat, betete 50 Jahre täglich für drei Freunde, dass sie Jesus kennenlernen. Zwei haben sich bekehrt als er schon alt war und der dritte erst nach seinem Tod. Kein Gebet ist vergebens \biblezit{IJoh}{5:14-15}. Aber was aus den Gebeten wird und auf welche Art sie ausgeführt werden, das weiss nur Gott.

Wenn wir uns die Gebete in der Bibel näher anschauen, so sind diese, immer ähnlich aufgebaut.
\begin{enumerate}
    \item Lobpreis Danksagung
    \item Busse
    \item Bitte
\end{enumerate}
Dieses Gebet von Nehemia in den Versen 5--11, ist so aufgebaut, wie das Vaterunser, welches wir aus \biblezit{Matt}{6:9} kennen.

\biblezit{Neh}{1:5-6}

In Vers 5 preist und lobt er Gott. In den Versen 6 und 7 tut Nehemia Busse. Nehemia betet aber nicht, dass die anderen sich vergangen haben und Fehler gemacht haben. Er identifiziert sich mit seinem Volk und spricht immer von \enquote{wir}. So am Schluss von Vers 6 \enquote{Auch \betonung[b, p]{ich} und das Haus meines Vaters haben gesündigt!} Das heisst nicht, dass wir für die Sünden unserer Vorfahren haften müssen, sondern nur, dass keiner besser als der andere ist und wir alle Sünder sind.

\biblezit{Neh}{1:7}
Auch so in Vers 7 \enquote{\betonung[b, p]{Wir} haben sehr verwerflich gegen dich gehandelt}. Nehemia war ja Mundschenk beim Persischen König, er hatte doch nichts mit den Verfehlungen seines Volkes zu tun. Stimmt, direkt nicht, aber er war Jude. Er wusste, wieso seine Vorfahren Israel verlassen mussten und anscheinend hat er auch die Bibel gelesen, weil er wusste, weshalb sie im Exil sind. Wir sagen ziemlich schnell, \enquote{Ach Gott warum lässt du das zu?}

Wenn wir sehen, wie die Menschen heute mit dem Gott, der sie erschaffen hat, umgehen, kann man sich nur wundern, dass die Welt nicht schon von ihm vernichtet wurde. Wenn Gott dieses Unrecht bestraft, sind auch wir mitten drin. Wir haben zwar die Sicherheit, dass wir nach dem Tod bei Jesus sind. Solange wir aber hier auf Erden sind, müssen wir trotzdem das auch mitmachen. Als die Israeliten wegen ihrer Sünden gegenüber Gott verschleppt wurden, wurden auch die gläubigen Juden nicht ausgeschlossen. Erst wenn wir entrückt werden, sind wir von der zukünftigen Trübsal befreit.

\biblezit{Neh}{1:8}

Ab Vers 8 erinnert Nehemia Gott an seine Zusagen. Hier ist nochmals ein Beweis, dass er die Bibel kannte.

Kennst \betonung[b, p]{du} dich in der Bibel so gut aus, dass du die Verheissungen Gottes kennst? Es ist wichtig, dass wir seine Zusagen und Verheissungen kennen. Wenn du die Bibel und die Verheissungen Gottes nicht kennst, kann es sein, dass du etwas erbittest, was Gott so nie versprochen hat und wunderst dich dann, dass du nicht erhört wirst und gibst dann Gott die Schuld. Zum Beispiel hat Gott nie versprochen, dass wir sofort gesund werden, oder Reich werden, oder Karriere machen, wenn wir darum beten. Gott kann uns oder einen lieben Angehörigen heilen und ja, das tut er heute noch, aber das Wann und Wie ist die souveräne Entscheidung Gottes. Wichtig! Es ist \betonung[b, p]{immer} Gott der heilt, nie ein Guru oder einer der die Hand auflegt.

Er hat uns aber versprochen, dass er uns in unseren Nöten unterstützt und trägt. Darauf kannst du dich berufen. Wie und wann er hilft, müssen wir ihm überlassen. 
\begin{quote}
    Zum Beispiel: \enquote{Herr, hast du in deinem Wort nicht versprochen, das du uns nur soviel auferlegst, wie ich tragen kann? Herr, du hast doch uns deine Unterstützung versprochen}     
\end{quote}
Nein, nur weil du jetzt das so betest, geht es nicht schneller. Gott weiss ja wie der Plan aussieht. So beruft sich auch Nehemia auf die Zusagen, die Gott seinem Volk in der Schrift gemacht hat.

\biblezit{Neh}{1:10-11}

In den Versen 10--11 bringt Nehemia dann sein Anliegen vor. Ihr Lieben, das ist der gefährlichste Teil des Gebets und sollte gut überlegt sein.

Es ist das eine, Gott um Hilfe zu bitten. Wie die Hilfe aussieht, das wissen wir nicht. Hätten wir nicht am liebsten, dass Gott mit den Fingern schnippt und als ist geregelt? Nachdem ich zum Glauben kam und ein bisschen verstand, worum es geht, habe ich gebetet, dass Gott doch auch in Naters eine FEG hinstellen möge. Es war in der Zeit, wo die FEG-Visp plante, wie es nun weiter gehen könnte. Da sie immer grösser wurden und es immer mehr Kinder gab, die betreut werden mussten, hatten sie ein paar Ideen wie sie das in Zukunft regeln wollen. Meine Idee war, dass sie sich nach Naters ausbreiten würden, um auch dort das Evangelium zu bringen. Aber das sollen die Pastoren der FEG machen, nicht ich.

Es kam aber anders als ich mir ausgedacht habe. Gott hatte andere Pläne als ich. Gott hat sich wohl gesagt, wenn du das willst, dann mach es doch. So, jetzt stehe ich hier, predige und darf in Gottes Reich mitarbeiten. Ich fühlte mich von Anfang an und noch jetzt, unfähig für diese Arbeit. Aber bis jetzt geht es hier vorwärts und ganz sicher nicht, weil ich so toll bin. 

Gott hat dich und mich an diesen Platz hier gestellt und hier und jetzt benutzt er uns. Uns zusammen als Team. Oft beten wir, Herr gebrauche mich und wenn es dann so weit ist, stöhnen wir. 

Es geht nie um uns. Wie Thomas das am Donnerstag so schön dargelegt hat. Wenn Jesus in \biblezit{Mk}{10:45} sagt, dass er zum Dienen hier ist, was sind denn wir? Wollen wir uns hier als Herren aufspielen? Ja, Gott gebraucht uns nicht immer so wie wir es gerne hätten, sondern so wie \betonung[b, p]{ER} es will und es geht \betonung[b, p]{immer um seine Verherrlichung}. So lesen wir in \biblezit{Jes}{42:16}, wie Gott führt. Oder wie es auch in Sprüche \biblezit{Spr}{20:24} steht. Ich weiss nicht, wenn Nehemia wirklich gewusst hätte, was auf ihn zukommt, ob er auch so intensiv gebetet hätte? Natürlich konnte er sich ausrechnen wie anstrengend und gefährlich es wird. Ganz sicher hatte er schon einen Plan im Kopf. Aber, das Resultat seiner ganzen Mühe? Darauf komme ich nochmals zurück

\section{Teil 3: Nehemia's Gebet wird erhört 2,1-8}
Wenn wir im Kapitel 2 weiter lesen, lesen wir, wie Gott das Gebet von Nehemia erhört hat. Es kam kein Blitz vom Himmel und die Mauern standen, es lief alles ziemlich weltlich ab.

\biblezit{Neh}{2:1}

Aber lesen wir den Vers 1 in Kapitel 2. Da steht zuerst eine neue Zeitangabe. Im Vers 1 Kapitel 1, heisst es im Monat Kislew im zwanzigsten Jahr Artaxerxes und hier heisst es im Monat Nisan, im zwanzigsten Jahr Artaxerxes. Kislew ist von Ende November bis Anfangs Dezember und Nisan ist nach unserer Rechnung von Mitte März bis Mitte April. Also 4 Monate nachdem er die Information von seinem Bruder bekam. Das Leben ging in der Zeit für Nehemia weiter. Er hat damals nicht alles hingeschmissen und ist abgereist. Er ging täglich seinen Pflichten nach, aber er betete täglich für Hilfe und eine Lösung für sein Problem. Vielleicht hat er auch schon gedacht, Gott hört mich nicht, oder vielleicht will er mich nicht hören? Aber er betete immer weiter, jeden Tag.

Jetzt, vier Monate nachdem er die Information über Jerusalem bekommen hatte, konnte er seine fröhliche Maske nicht mehr aufrechterhalten und der König sah ihm an, dass mit ihm etwas nicht stimmte. Wenn wir betrübt sind, hat das nicht nur psychische Auswirkungen, sondern auch körperliche. Darum sollten wir uns gegenseitig beobachten und nachfragen, wenn wir Veränderungen an unseren Mitgeschwistern wahrnehmen. Oft zeigen sich Probleme darin, dass sich die Person zurückzieht und die Gesellschaft meidet.

Jetzt, in dieser Situation sehen wir, wie vertraut der Mundschenk mit dem König war. Ob jetzt der König wirklich aus Mitgefühl fragte, weil er sich Sorgen um Nehemia machte, oder eher wegen sich selber, weiss ich nicht. Aber es ist ihm aufgefallen, dass Nehemia traurig aussah. Das Fasten und Beten und die Sorgen um sein Volk haben seine Spuren hinterlassen. Hier zeigt es sich, dass Gott auch die Könige nach seinem Willen führt und leitet wie wir in \biblezit{Spr}{21:1} lesen. Auch Weltherrscher sind nur Knechte Gottes, auch wenn sie denken, sie seien die Macher.

In jener Zeit in der Nähe eines Machthabers zu arbeiten, war nicht nur nicht einfach, sondern auch gefährlich. Man war ihren Launen ausgesetzt. Auch bei Esther lesen wir, welche Angst sie hatte vor den König zu treten, \biblezit{Est}{4:11} so auch hier, Nehemia. Darum sagte er: \enquote{Da fürchtete ich mich sehr}.

Nehemia fasste sich ein Herz und erklärte dem König die Situation. Ich denke, er hat nicht erwartet, dass der König ihm helfen würde. Als dann die Antwort vom König war: \enquote{Was erbittest du denn?} War er sicher sehr überrascht.

Seine Reaktion darauf war als Erstes ein Stossgebet zu Gott. Passiert dir das auch in Situationen, wo du überrascht wirst? Wir haben hier im Wallis so ausdrücke wie: \betonung[b, p]{O, Gott} oder \betonung[b, p]{Jessus Gott}, das sind alles Stossgebete, die zu Gott gehen. Das sollte uns immer bewusst sein. Oft passiert es aus Gewohnheit, aber noch schlimmer ist es, wenn wir es im Zorn sagen. So sagt Jakobus in seinem Brief \biblezit{Jak}{3:9}? Gottes Namen zu missbrauchen ist nicht harmlos. Gott hat es in den 10 Geboten im 3. Gebot \biblezit{IIMos}{20:7} gesagt, dass wir seinen Namen nicht missbrauchen sollen.

\biblezit{Neh}{2:5-6}

Wenn der König fragt, musste man Antworten. In dieser Situation hatte man nicht mehr lange Zeit zu überlegen. Und einem König nicht zu Antworten war tödlich. Aber Nehemia war vorbereitet.

Ich denke nicht, dass der König Mitleid mit den Menschen in Israel und Jerusalem hatte. Wobei, Nehemia bei seiner Schilderung Jerusalem gar nicht erwähnte. Wohl aus weiser Vorsicht, dass dies den König verärgern könnte. Vielleicht war der König beeindruckt über die Loyalität, die Nehemia gegenüber seinem Volk zeigte. Was sicher ist, ist, dass Gott die ganze Situation im Griff hatte und leitete. 

Auch wenn Gott den König führte, Nehemia musste selber vor dem König stehen und seine Anliegen vorbringen. Auch wir sollten nach dem Gebet Taten folgen lassen, wenn sich die Gelegenheit bietet. Nehemia hat nicht direkt, nachdem er die schreckliche Nachricht von seinem Bruder gehört hat, eigene Aktionen gestartet. Er hat so lange gebetet, bis sich ihm eine Möglichkeit auftat, einen Schritt weiter zu gehen. Dann als es so weit war, hat er die Gelegenheit beim Schopf gepackt.

Das sollten wir beachten. Wenn wir eine Idee haben, sollen wir nicht direkt draufloslaufen, sondern zuerst beten, planen und warten bis Gott eine Gelegenheit schafft. Vielleicht fragst du dich: \enquote{Wie merke ich denn, dass die Zeit zum Handeln gekommen ist?} Meine Antwort; ich weiss es nicht. Das ist individuell verschieden. Ich weiss nur, dass Gott es ganz sicher so einrichten wird, dass du es merken wirst, dass die Zeit da ist. Wie bei unserem Adventsmarkt. Wir haben die Unterlagen abgegeben, gebetet und gewartet. Wider Erwarten kam ein positiver Entscheid. Jetzt mussten wir uns aber nochmals entscheiden, hätten immer noch Nein sagen können. Ich denke, dann hätten wir gegen Gottes Willen entschieden. Jetzt müssen wir an die Arbeit. Nicht überstürzt, sondern überlegt mit reichlich gemeinsamen Gebet. Darum ist auch der Gebetsabend so wichtig.

\biblezit{Neh}{2:7-8}
Nehemia war vorbereitet. Nehemia hatte seinen Plan schon Kopf. Vielleicht hatte er sich in den letzten vier Monaten bei Bauexperten erkundigt, wie man so eine Mauer repariert. Er hat dem König bei seiner Nachfrage sogar eine Zeit angegeben können, wie lange er weg sein würde. In \biblezit{Neh}{5:14} lesen wir, dass er 12 Jahre Statthalter in Jerusalem war, dann wieder zurückkehrte, um später nochmals nach Jerusalem zurückzukehren.

Nehemia dachte auch an Details wie Empfehlungsschreiben für unterwegs. Organisierte das Baumaterial für den Aufbau. Auch das Reisen war nicht so einfach wie heute. Sich in einen Flieger setzen und nach Jerusalem fliegen das ging nicht. Nehemia musste eine Karawane zusammenstellen. Er brauchte Kamele, Kameltreiber,  Krieger zum Schutz der Karawane, Reiseführer, Proviant für mehrere Monate. Es waren ca. 1700 km. Wenn eine Karawane ca. 22 - 25 km pro Tag zurücklegte, so dauerte die Reise ca. 3 Monate. Durch Wüsten, Gebirge, Kälte und Hitze.
Sein Vorhaben war also nicht, ein kleiner Weihnachtsmarkt. Es wäre einfacher gewesen in der Burg zu bleiben und zu Gott zu beten und zu hoffen, dass sich jemand anderes die Menschen in Jerusalem kümmert.

\section{Schluss}

\biblezit{Neh}{2:8}

Wir haben nun gesehen wie Nehemia Informationen sammelte, für Hilfe betete, wie er durch den König Hilfe bekam, wie er diese Hilfe annahm und dann Nägel mit Köpfen machte. Sein Ausspruch nach alldem war \biblezit{Neh}{2:8}b \betonung[b, p]{weil die gute Hand meines Gottes über mir war.}

Dieser Text steht unter anderem in der Bibel, um uns Mut zu machen geduldig zu Beten und wenn sich die Gelegenheit ergibt, diese beim Schopf zu packen. Gott wird einem das nötige Rüstzeug dazu geben, sodass es auch funktionieren wird. Wir sollten nicht überstürzt handeln, sondern immer im Gebet auf Gott harren. So das auch du, sagen kannst, \betonung[b, p]{weil die gute Hand meines Gottes über mir war.}

Wir sind menschlich gesehen, hier nur eine kleine Gruppe, aber wir haben einen mächtigen Chef, der sich nie irrt, der nie einen Fehler macht, der alles kann und alles weiss. Auf ihn können wir uns verlassen.

Das wusste auch Paulus, darum schreibt er in \biblezit{Rom}{8:31} folgendes. Und wisst ihr was das Beste ist? Wenn wir uns an die Regeln Gottes halten, seine Gösse im Blick haben und Schritt für Schritt vorwärtsgehen, können wir nichts falsch machen. Egal wie das Resultat ausfällt. Paulus wurde hingerichtet und die Gemeinden gingen sehr schnell eigene Wege. Irrlehrer tauchten auf, Religionen wurden erfunden, neue Gesetze und Rituale ins Leben gerufen. Paulus hat seinen Kampf gekämpft und wird dafür im Himmel belohnt. Darum lasst uns unseren Kampf kämpfen. Den Rest überlassen wir Gott. Den Lohn dürfen wir erst im Himmel erwarten.

Nehemia's Geschichte ist hier nicht zu Ende. Ihr könnt zu Hause die Geschichte zu Hause zu Ende lesen. Nehemia hatte mit vielen Widerständen und Widrigkeiten zu kämpfen. Immer wieder ging er ins Gebet und bittet Gott um seine Hilfe. Lasst uns beten!

\end{document}


